\section{プログラミングの基礎}
\subsection{授業内容}
\subsubsection{科目の中でのこのコマの位置づけ}
後期の授業の主眼は「データから意味のある情報を取り出す」ことにある。すなわち，これまでのデータ駆動型(Data Driven)な発想を逆転し，データがどのように生成されたと考えるかという観点から，データ生成モデル(Data generating Model)による理解を試みる。

モデルは数学的表現がなされ，モデルの形成やデータとの照合(fitting)には計算機の利用が必須である。後期の初回となる今回の目的は，プログラミングの基本的な考え方を身につけ，次回以降の本格的な運用に備えることである。
プログラミングの基礎はコマンドによる命令であり，コマンドの書き間違えはエラーとなって帰ってくる。一見不親切に思えるが，即時反応により学習の効率は良く，コード補完機能などを用いることで簡単なミススペルは回避することができる。小さなものから大きくしていくこと，1行ずつ実行することなど，基本的な姿勢についても理解する。

またプログラミング言語(高級言語)はいくつかあるが，文法的な違いを除けばその本質は代入，反復，条件分岐である。この三点について理解し，基本的な書き方を学ぶ。実際にRでコードを書きながら，その挙動について確認する。
最終的な到達段階として，Fizz-Buzz問題や行列計算ができるコードを書くこととする。

\subsubsection{コマ主題細目}
\begin{description}
    \item[プログラミングの基礎] プログラミングにあたって重要なことは「思った通りに」動くのではなく，「書いた通りに」動くことである。ミススペルや大文字小文字の違いにも注意が必要である。コードのスペルチェックや補完機能を活用し，また1行ずつ確認しながら進めるといったプログラミングに関わる基本的な心構えについて理解する。

    \item[いくつかのプログラミング言語] Rはプログラミング言語の一種であると言っても良い。プログラミング言語には他にもPythonやBasic，C言語などがある。これらの基本的な関係について理解するとともに，コンパイラとインタプリタという実行形式の違いについても理解する。この違いは今後確率的プログラミング言語を利用する際の知識として生きてくる。

    \item[高級言語の基本的な働き]高級言語と呼ばれるプログラミング言語の基本的な働きは，代入，反復，条件分岐である。Rでは\verb|<-|や\verb|=|で代入を，\verb|for|や\verb|while|で反復を，\verb|if|や\verb|if_else|で条件分岐を行う。これらの表現は言語間を通じて共通なものが多いため，その基本的な振る舞いを確認することは技術の一般化に役立つ。

    \item[入れ子になったインデックス] 代入と反復を組み合わせることで行列の計算など煩雑なものが文字通り機械的に計算されていくことになる。さらに「行の番号」を代入することで個体を識別し，それがさらに別の変数のインデックスとなる入れ子になった計算についても練習しておくことで，今後の階層モデルへの準備になる。

\end{description}
\subsubsection{キーワード}
\begin{itemize}
    \item プログラミング言語
    \item コンパイル
    \item 条件分岐
    \item 階層性
\end{itemize}
\subsection{授業情報}
\paragraph{コマの展開方法}
講義/遠隔/演習
\subsubsection{予習・復習課題}
\paragraph{予習}事前に環境の準備をしておく必要がある。環境の準備についてはいくつかの方策があり，これについては導入資料を参照しながら準備しておくこと。なお，環境準備中に問題が生じた場合はいち早く教員かTAに相談し，実行できるようにしておくこと。
\paragraph{復習}反復計算の練習課題，条件分岐の練習課題など，複数の課題にしっかり取り組むこと。